\subsection{Koordinatendarstellung von Vektoren}
\label{subsec:vector_algebra_coordinate_representation}

In einem Koordinatensystem wird eine Verschiebung durch ihre Anteile entlang
der Koordinatenachsen beschrieben.

\begin{DefinitionBox}{Koordinaten eines Vektors}{vector_algebra_vector_coordinates}
  Ein Vektor im $n$-dimensionalen reellen Raum besitzt die
  \emph{Koordinatendarstellung} \EnglishTerm{coordinate representation}
  \[
    \vec a=
    \begin{pmatrix}
      a_1\\a_2\\\vdots\\a_n
    \end{pmatrix}
    \in\R^n.
  \]
  Die Zahl $a_i$ gibt den gerichteten Anteil von $\vec a$ entlang der
  $i$-ten Koordinatenachse an.
\end{DefinitionBox}

\noindent
\begin{minipage}[c]{0.48\linewidth}
  \[
    \vec a=\begin{pmatrix}1\\2\end{pmatrix},
    \qquad
    \vec b=\begin{pmatrix}-2\\1\end{pmatrix}
    \quad\in\R^2.
  \]
  Die Pfeile dürfen parallel verschoben werden; ihre Koordinaten bleiben
  unverändert.
\end{minipage}\hfill
\begin{minipage}[c]{0.48\linewidth}
  \centering
  \CachedTikzFigure{content/resources/chapters/vector_algebra/tikz/coordinate_vectors_2d.tex}
\end{minipage}

\begin{ExampleBox}{Vektorkoordinaten ablesen}{vector_algebra_read_coordinates}
  Die gestrichelten Hilfslinien zerlegen die Verschiebung entlang der Achsen:
  \[
    \vec a=
    \begin{pmatrix}
      \Cloze{1.2cm}{3}\\[0.35em]
      \Cloze{1.2cm}{2}
    \end{pmatrix}
    \in\R^2.
  \]
  \begin{center}
    \CachedTikzFigure{content/resources/chapters/vector_algebra/tikz/read_vector_coordinates.tex}
  \end{center}
\end{ExampleBox}

\noindent
\begin{minipage}[c]{0.43\linewidth}
  Im dreidimensionalen Raum kommt die Verschiebung entlang der $x_3$-Achse
  hinzu. Beispielsweise gilt
  \[
    \vec c=\begin{pmatrix}2\\1\\3\end{pmatrix}\in\R^3.
  \]
\end{minipage}\hfill
\begin{minipage}[c]{0.53\linewidth}
  \centering
  \CachedTikzFigure{content/resources/chapters/vector_algebra/tikz/coordinate_vector_3d.tex}
\end{minipage}

\begin{DefinitionBox}{Nullvektor}{vector_algebra_zero_vector}
  Der \emph{Nullvektor} \EnglishTerm{zero vector} besitzt ausschließlich
  Nulleinträge:
  \[
    \vec 0=
    \begin{pmatrix}0\\\vdots\\0\end{pmatrix}\in\R^n,
    \qquad
    \vec 0=\begin{pmatrix}0\\0\end{pmatrix}\in\R^2,
    \qquad
    \vec 0=\begin{pmatrix}0\\0\\0\end{pmatrix}\in\R^3.
  \]
  Er beschreibt eine Verschiebung der Länge null und besitzt keine Richtung.
\end{DefinitionBox}
